% BEGIN VOL07-EXPANSION VII42
\section{Supplementary worked examples}
\label{sec:vii42-pedagogy}

\begin{example}[Cylinder line fields]\label{ex:vii42-ped-01}
On a circular cylinder of radius \(R\), principal curvatures are \(1/R\) and \(0\) up to orientation.  The principal directions follow circles and generators, so curvature lines form two orthogonal families with no umbilics.
\end{example}

\begin{example}[Sphere umbilics]\label{ex:vii42-ped-02}
On a round sphere, the shape operator is a scalar multiple of the identity at every point.  Every tangent direction is principal, so no distinguished principal line field exists anywhere; the entire sphere is umbilic.
\end{example}

\begin{example}[Monge-patch estimate]\label{ex:vii42-ped-03}
For \(z=\tfrac12(ax^2+2bxy+cy^2)\) at a horizontal tangent point, the Hessian \(\begin{pmatrix}a&b\\b&c\end{pmatrix}\) approximates the shape operator.  Its eigenpairs provide local principal-curvature and direction estimates used in mesh fitting.
\end{example}

\section{Graded supplementary exercises}
The following exercises add a deliberate progression from concrete calculation to structural proof, hypothesis testing, application, and synthesis.

\subsection*{Standard computations and constructions}

\begin{exercise}[Cylinder curvatures]\label{exr:vii42-ped-01}
State the principal curvatures of a radius-\(R\) cylinder up to normal orientation.
\end{exercise}

\begin{hint}
One direction bends, one does not.
\end{hint}

\begin{solution}
\(1/R\) and \(0\), with simultaneous sign reversal if the normal is reversed.
\end{solution}

\begin{exercise}[Sphere umbilic]\label{exr:vii42-ped-02}
Why is every point of a round sphere umbilic?
\end{exercise}

\begin{hint}
Its shape operator is scalar.
\end{hint}

\begin{solution}
Both principal curvatures are equal at every point.
\end{solution}

\begin{exercise}[Euler direction]\label{exr:vii42-ped-03}
What normal curvature is obtained when \(\theta=0\) in Euler's formula?
\end{exercise}

\begin{hint}
Align with \(e_1\).
\end{hint}

\begin{solution}
\(k_1\).
\end{solution}

\begin{exercise}[Euler quarter turn]\label{exr:vii42-ped-04}
What normal curvature is obtained at \(\theta=\pi/2\)?
\end{exercise}

\begin{hint}
Align with \(e_2\).
\end{hint}

\begin{solution}
\(k_2\).
\end{solution}

\begin{exercise}[Hessian eigenvalues]\label{exr:vii42-ped-05}
For \(\begin{pmatrix}2&0\\0&-1\end{pmatrix}\), what are the small-slope principal-curvature estimates?
\end{exercise}

\begin{hint}
Read the diagonal eigenvalues.
\end{hint}

\begin{solution}
\(2\) and \(-1\).
\end{solution}

\subsection*{Proofs}

\begin{exercise}[Orthogonal families]\label{exr:vii42-ped-06}
Why do the two curvature-line families cross orthogonally away from umbilics?
\end{exercise}

\begin{hint}
The shape operator is self-adjoint.
\end{hint}

\begin{solution}
Eigenvectors for distinct eigenvalues of a self-adjoint operator are orthogonal.
\end{solution}

\begin{exercise}[Line not vector]\label{exr:vii42-ped-07}
Why must a numerical principal-direction field identify \(e\) with \(-e\)?
\end{exercise}

\begin{hint}
Eigenvectors have sign ambiguity.
\end{hint}

\begin{solution}
Both vectors span the same eigendirection, so the geometric object is a line field.
\end{solution}

\begin{exercise}[Ridge first derivative]\label{exr:vii42-ped-08}
What first-order condition defines a \(k_1\)-ridge candidate under the chapter convention?
\end{exercise}

\begin{hint}
Differentiate \(k_1\) along \(e_1\).
\end{hint}

\begin{solution}
\(D_{e_1}k_1=0\).
\end{solution}

\begin{exercise}[Ridge second derivative]\label{exr:vii42-ped-09}
What additional sign selects a strict \(k_1\)-ridge?
\end{exercise}

\begin{hint}
It is a local maximum along \(e_1\).
\end{hint}

\begin{solution}
\(D_{e_1}^2k_1<0\).
\end{solution}

\subsection*{Counterexamples and hypothesis tests}

\begin{exercise}[Ridge equals curvature line]\label{exr:vii42-ped-10}
Test the claim.
\end{exercise}

\begin{hint}
One is an implicit derivative-zero locus, the other an integral curve.
\end{hint}

\begin{solution}
False.  They coincide only under additional geometric conditions.
\end{solution}

\begin{exercise}[Principal direction at umbilic unique]\label{exr:vii42-ped-11}
Test the claim.
\end{exercise}

\begin{hint}
At an umbilic the shape operator is scalar.
\end{hint}

\begin{solution}
False.  Every direction is principal, so there is no distinguished eigenline.
\end{solution}

\begin{exercise}[One-ring always best]\label{exr:vii42-ped-12}
Test the claim for noisy mesh curvature estimation.
\end{exercise}

\begin{hint}
Curvature is a second-derivative quantity.
\end{hint}

\begin{solution}
False.  A one-ring can be too noisy; scale selection and robust fitting are essential.
\end{solution}

\subsection*{Applications and investigations}

\begin{exercise}[Feature extraction]\label{exr:vii42-ped-13}
Why are ridge and valley curves useful in geometry processing?
\end{exercise}

\begin{hint}
They track extrema of principal curvature.
\end{hint}

\begin{solution}
They highlight perceptually and mechanically salient shape features useful for segmentation, stylization, inspection, and matching.
\end{solution}

\begin{exercise}[Scale-space]\label{exr:vii42-ped-14}
Why compare ridge/valley features over several neighborhood scales?
\end{exercise}

\begin{hint}
Small scales amplify noise; large scales erase detail.
\end{hint}

\begin{solution}
Persistence across scales helps separate stable geometric structure from discretization noise.
\end{solution}

\subsection*{Challenge problems}

\begin{exercise}[Monge eigenvectors]\label{exr:vii42-ped-15}
For \(H=\begin{pmatrix}3&0\\0&1\end{pmatrix}\), identify principal directions in the fitted tangent frame.
\end{exercise}

\begin{hint}
Use matrix eigenvectors.
\end{hint}

\begin{solution}
They are the coordinate axes, with estimated principal curvatures \(3\) and \(1\).
\end{solution}

\begin{exercise}[Sign-consistency challenge]\label{exr:vii42-ped-16}
Why can naively interpolating principal eigenvectors across a mesh create artificial flips?
\end{exercise}

\begin{hint}
Each eigenvector may independently change sign.
\end{hint}

\begin{solution}
Because the field is unoriented, neighboring estimates can differ by sign; algorithms must compare lines or explicitly choose consistent signs before interpolation/tracing.
\end{solution}

% END VOL07-EXPANSION VII42
